Entscheiden Sie, ob die folgenden Funktionen eine elementare Stammfunktion
haben:
\begin{teilaufgaben}
\item $f(x)= (1+\log(x))x^x$
\item $f(x)=x^x$
\end{teilaufgaben}

\begin{loesung}
Wir schreiben $x^x = e^{x \log(x)}$, die Funktion
$\vartheta_2 = x^x = \exp(x\log(x))$ ist also eine exponentielle Erweiterung
des Körpers $K(x,\vartheta_1)$ mit $\vartheta_1=\log x$.
$\vartheta_2$ hat die definierende Eigenschaft
$D\vartheta_2 = \vartheta_2 \cdot x\log(x)$.
\begin{teilaufgaben}
\item 
Für dieses Integral wird die Risch-Differentialgleichung zu
\[
Dq_1 + (\log(x)+1)q_1 = 1 + \log(x),
\]
die die Lösung $q_1=1$ hat.
Es folgt, dass
\[
\int (1+\log(x))x^x = 1\cdot \vartheta_2 = x^x+C.
\]
Tatsächlich ist die Ableitung
\begin{align*}
Dx^x
&=
D(\exp(x\log(x))
=
\exp(x\log(x)) \biggl(\log(x) + x\cdot\frac{1}{x}\biggr)
=
(1+\log(x))x^x.
\end{align*}
\item
Da $\vartheta_2$ eine exponentielle Erweiterung ist, muss die
Risch-Differentialgleichung
\[
Dq_1
+
Du\, q_1 = p_1
\qquad\Rightarrow\qquad
Dq_1 + (\log(x) + 1)q_1 = 1
\]
mit $q_1\in K(x,\vartheta_1)$ gelöst werden.
Eine Lösung kann als gekürzter Bruch $a(\vartheta_1)/b(\vartheta_1)$ mit
Polynomen $a(\vartheta_1),b(\vartheta_1)\in K(x)[\vartheta_1]$ geschrieben
werden, wobei $b(\vartheta_1)$ ausserdem normiert sein soll.
Einsetzen in die Differentialgleichung und multiplizieren mit $b(\vartheta_1)^2$
ergibt
\[
Da(\vartheta_1) b(\vartheta_1)
-
a(\vartheta_1)Db(\vartheta_1)
+
(\vartheta_1+1)a(\vartheta_1)b(\vartheta_1)
=
b(\vartheta_1)^2.
\]
Da alle Terme ausser möglicherweise dem zweiten durch $b(\vartheta_1)$
teilbar sind, muss auch der zweite teilbar sein.
Weil der Bruch gekürzt ist, muss $b(\vartheta_1)\mid Db(\vartheta_1)$
sein.
Da aber $Db(\vartheta_1)$ kleineren Grad als $b(\vartheta_1)$ hat,
muss $Db(\vartheta_1)=0$ und damit $b(\vartheta_1)=b_0$ eine Konstante sein.
Da $b(\vartheta_1)$ normiert ist, folgt ausserdem $b_0=1$.

Die Differentialgleichung vereinfacht sich jetzt zu
\[
\begin{aligned}
&&
Da(\vartheta_1) + (\vartheta_1+1) a(\vartheta_1)
&= a(\vartheta_1)
\\
&\Rightarrow&
Da(\vartheta_1)
&= -\vartheta_1 a(\vartheta).
\end{aligned}
\]
Auf der rechten Seite steht ein Polynom vom Grad $1+\deg a(\vartheta)$
während die linke Seite höchstens Grad $\deg a(\vartheta_1)$ hat.
Somit kann es keine Lösung der Differentialgleichung geben, das Integral
ist nicht elementar.
\qedhere
\end{teilaufgaben}
\end{loesung}
